\documentclass[12pt, twoside]{article}
%\documentclass[12pt, twoside]{article}
\usepackage[letterpaper, margin=1in, headsep=0.2in]{geometry}
\setlength{\headheight}{0.6in}
%\usepackage[english]{babel}
\usepackage[utf8]{inputenc}
\usepackage{microtype}
\usepackage{amsmath}
\usepackage{amssymb}
%\usepackage{amsfonts}
\usepackage[nomessages]{fp} %\FPeval{\var-name}{2*sin(pi/6)}
\usepackage{siunitx} %units in math. eg 20\milli\meter
\usepackage{yhmath} % for arcs, overparenth command
\usepackage{tikz} %graphics
\usetikzlibrary{quotes, angles, arrows, arrows.meta}
\usepackage{pgfplots}
\usepgfplotslibrary{fillbetween}
\pgfplotsset{compat=1.16}
\usepackage{graphicx} %consider setting \graphicspath{{images/}}
\usepackage{parskip} %no paragraph indent
\usepackage{enumitem}
\usepackage{multicol}
\usepackage{venndiagram}

\usepackage{fancyhdr}
\pagestyle{fancy}
\fancyhf{}
\renewcommand{\headrulewidth}{0pt} % disable the underline of the header
\raggedbottom
\hfuzz=2mm %suppresses overfull box warnings

\usepackage{hyperref}

\fancyhead[LO]{La Scuola d'Italia / Dr.\ Huson / IB Math 11 \\* 27 October 2026}
\fancyhead[RO]{\fbox{\begin{minipage}{6cm}\footnotesize Nickname:\\[0.5cm]\end{minipage}}}
\fancyfoot[C]{\thepage}

\begin{document}
\subsubsection*{2.5 Pre-Quiz: Sequences, Lines, Exponents and Radicals}

A calculator is allowed. When a question asks for an exact answer or a
particular form, show the algebra that leads to it.

\begin{enumerate}[itemsep=0.15cm]

%--- Problem 1: Arithmetic sequence from two terms: d and u_1, u_n, S_20, show 199 is a term and which term [10 marks] ---
%--- Source: Original (AA SL 1.2); review of unit 1, cf. 1-10Q_Sequences P1-P2 ---
\item[1.]
In an arithmetic sequence, $u_{4} = 23$ and $u_{11} = 51$.

\begin{enumerate}[itemsep=0.1cm]
  \item Find the common difference $d$ and the first term $u_{1}$.
  \item Find an expression for $u_{n}$ in terms of $n$.
  \item Find $S_{20}$, the sum of the first $20$ terms.
  \item Show that $199$ is a term of the sequence, and find which term
  it is.
\end{enumerate}

\begin{tikzpicture}
  \draw (-0.6,0) rectangle (14.6,15.3);
  \foreach \y in {1,2,...,14} \draw[dotted] (0,\y)--(14.1,\y);
\end{tikzpicture}

\newpage

%--- Problem 2: Geometric sequence from two terms: r and u_1, S_8, least n with S_n > 100 000 [8 marks] ---
%--- Source: Original (AA SL 1.3); review of unit 1, cf. 1-10Q_Sequences P3 ---
\item[2.]
In a geometric sequence, $u_{2} = 12$ and $u_{5} = 324$.

\begin{enumerate}[itemsep=0.1cm]
  \item Find the common ratio $r$ and the first term $u_{1}$.
  \item Find $S_{8}$, the sum of the first $8$ terms.
  \item Find the least value of $n$ for which $S_{n} > 100\,000$.
\end{enumerate}

\begin{tikzpicture}
  \draw (-0.6,0) rectangle (14.6,18.2);
  \foreach \y in {1,2,...,17} \draw[dotted] (0,\y)--(14.1,\y);
\end{tikzpicture}

\newpage

%--- Problem 3: Gradient of PQ; perpendicular L through R; foot F (not the midpoint); RF in surd form [11 marks] ---
%--- Source: Original, in the style of AA SL Paper 1 coordinate-geometry questions (AA SL 2.1, 3.1) ---
\item[3.]
The points $P(-4,\ 8)$, $Q(6,\ 3)$ and $R(4,\ 9)$ are shown on a
coordinate plane. The line $L$ passes through $R$ and is perpendicular
to the line $PQ$. The line $L$ meets the line $PQ$ at the point $F$.

\begin{enumerate}[itemsep=0.1cm]
  \item Find the gradient of the line $PQ$.
  \item Find the equation of $L$. Give your answer in the form
  $y = mx + c$.
  \item Find the coordinates of $F$.
  \item Find the distance $RF$. Give your answer in the form $a\sqrt{5}$,
  where $a \in \mathbb{Z}$.
\end{enumerate}

\begin{tikzpicture}
  \draw (-0.6,0) rectangle (14.6,16.3);
  \foreach \y in {1,2,...,15} \draw[dotted] (0,\y)--(14.1,\y);
\end{tikzpicture}

\newpage

%--- Problem 4: Museum tickets: set up and solve a 2x2 system; number of children in a group [7 marks] ---
%--- Source: Original (AA SL 2.1, systems of linear equations); in the style of IB linear-model questions ---
\item[4.]
At a museum, an adult ticket costs $x$ euros and a child ticket costs
$y$ euros. A group of $3$ adults and $2$ children pays $61$ euros. A group
of $2$ adults and $5$ children pays $70$ euros.

\begin{enumerate}[itemsep=0.1cm]
  \item Write down two equations in $x$ and $y$.
  \item Find the price of an adult ticket and the price of a child ticket.
  \item A school group of $4$ adults and some children pays $156$ euros.
  Find the number of children in the group.
\end{enumerate}

\begin{tikzpicture}
  \draw (-0.6,0) rectangle (14.6,17.2);
  \foreach \y in {1,2,...,16} \draw[dotted] (0,\y)--(14.1,\y);
\end{tikzpicture}

\newpage

%--- Problem 5: Exponent laws; exact fractional negative power; same-base equation; equation in u = 3^x [12 marks] ---
%--- Source: Original, in the style of AA SL Paper 1 exponent questions (AA SL 1.5, 1.7, 2.10) ---
\item[5.]

\begin{enumerate}[itemsep=0.1cm]
  \item Simplify $\left(3a^{2}b^{-1}\right)^{2} \times 2a^{-3}b^{4}$.
  \item Find the exact value of $\left(\dfrac{16}{81}\right)^{-\frac{3}{4}}$.
  \item Solve the equation $4^{2x - 1} = 8^{x + 2}$.
  \item By letting $u = 3^{x}$, solve the equation
  $3^{2x} - 10\left(3^{x}\right) + 9 = 0$.
\end{enumerate}

\begin{tikzpicture}
  \draw (-0.6,0) rectangle (14.6,17.5);
  \foreach \y in {1,2,...,16} \draw[dotted] (0,\y)--(14.1,\y);
\end{tikzpicture}

\newpage

%--- Problem 6: Surds: collect like surds, rationalize with a two-surd conjugate, linear equation with a surd coefficient, conjugate-product area [10 marks] ---
%--- Source: Original, in the style of AA SL surd questions (AA SL 1.7 / prior learning) ---
\item[6.]
In this question, give every answer in exact form.

\begin{enumerate}[itemsep=0.1cm]
  \item Simplify $\sqrt{98} - \sqrt{32} + \sqrt{18}$.
  \item Write $\dfrac{10}{\sqrt{7} - \sqrt{2}}$ in the form
  $a\sqrt{7} + b\sqrt{2}$, where $a, b \in \mathbb{Z}$.
  \item Solve the equation $x\sqrt{3} = x + 2$. Give your answer in the
  form $a + b\sqrt{3}$, where $a, b \in \mathbb{Z}$.
  \item A rectangle has sides of length $\left(4 + \sqrt{5}\right)\,\text{cm}$
  and $\left(4 - \sqrt{5}\right)\,\text{cm}$. Find the area of the
  rectangle.
\end{enumerate}

\begin{tikzpicture}
  \draw (-0.6,0) rectangle (14.6,15.8);
  \foreach \y in {1,2,...,15} \draw[dotted] (0,\y)--(14.1,\y);
\end{tikzpicture}

\end{enumerate}
\end{document}
